% =====================================================================
% Nuovo paragrafo per §V-D, dopo il capoverso su prior recalibration.
% Numeri: i3lung_prior_estimability.csv (30 transfer, 10 semi).
% =====================================================================

The reason prior adjustment behaves this way is measurable. The method needs the
target class prior, which it does not observe, so it estimates it from the
source model's predictions on unlabeled target covariates. We read that estimate
out and compared it against the truth on all thirty transfers. It is off by
$0.163$ in median absolute error, on a quantity whose true value spans $0.227$
across the cohort, and the failures are large: for four transfers the estimated
prevalence sits above $0.98$ where the true value is below $0.80$, and for four
others it falls below $0.35$ where the true value is above $0.64$. The damage
the method does follows that error closely (Spearman $\rho=-0.80$): where the
prior is estimated well the correction is harmless or helpful, and where it is
estimated badly the correction degrades calibration in proportion.

This locates the failure. The assumption prior adjustment makes is reasonable
and the correction it applies is correct given the prior; what breaks is the
estimate of a quantity that the target cohort does not reveal. It also explains
the heterogeneity reported above, since an effect governed by an estimation
error that ranges from $0.003$ to $0.574$ has no meaningful average.

A deployment cannot use this. Detecting that the estimate went wrong requires
comparing it against the target labels, which is the information the estimate
exists to replace, and inspecting the estimate for implausibility does not work
either: the transfer where prior adjustment performs best returns an estimated
prevalence of $0.982$, at the boundary of the parameter space. The thirty
transfers share six centers and are not independent, and this analysis was
carried out after the main results, so we report it as an explanation of an
observed pattern rather than as a confirmatory finding.
